\documentclass[11pt]{article}
\usepackage[margin=1in]{geometry}
\usepackage{amsmath,amssymb,amsthm}
\usepackage[T1]{fontenc}
\usepackage[hidelinks]{hyperref}
\newtheorem{proposition}{Proposition}
\newtheorem{lemma}[proposition]{Lemma}
\newtheorem{corollary}[proposition]{Corollary}
\newcommand{\F}{\mathbb F}
\newcommand{\C}{\mathbb C}
\newcommand{\A}{\mathbb A}
\newcommand{\ind}{\operatorname{ind}}
\newcommand{\SL}{\operatorname{SL}}
\title{Background search: global topological coupling and exact scope gates}
\author{Independent internal research audit}
\date{16 November 2026}
\begin{document}
\maketitle

\noindent\textbf{Status.}
This note records an independent background literature search and elementary
deductions. It does not close either full exceptional Mathieu case or the
Keller empty-fibre obstruction, and makes no novelty or priority claim.
No agents were messaged. Existing canonical proofs and registry entries were
not changed. The proofs below are written arguments, with small exact controls;
they are not specialist-reviewed or formally verified.

\section{Current progress and the remaining obligations}

The registry entry \texttt{GMCSOURCEHOMGEN20261116} records a general pointwise
Gaussian decision producer over the rationals and simple number fields, in
arbitrary dimension and degree. It retains all original markers, exact
integration-by-parts identities, finite actual Wick jets and sufficient
numerical cutoffs. A safe result and a fixed bad-marker certificate are both
possible outputs. This is an algorithmic result, not a positive universal
Gaussian conjecture. The current Gaussian thread is pursuing
specialization-valid coefficient-family machinery.

The canonical Mathieu regular-period saturation note
identifies the same unit-adjoined regular-period sector in every smooth proper
phase model, for every original phase. In its notation,
\[
 R_f=\C(s)u\oplus I_f,\qquad
 D_f=M_f/R_f,\qquad
 \epsilon_f\in\operatorname{Ext}^1_D(D_f,\C(s)).
\]
Under eventual purity, the missing assertion is $\epsilon_f=0$ on the
\emph{entire} exceptional pure locus, or an actual pure phase with nonzero
$\epsilon_f$ and a fixed original bad marker. Choosing a different proper
model cannot supply missing marked directions. This is a precise reduction,
not evidence that the remaining theorem is small.

A later update, \texttt{LOGARITHMIC\_JETS\_AND\_SPARSE\_ALL\_RING\_SAFETY.md},
closes MZ for every phase supported on at most four left-torus weights
on $\SL_2/2O$, and at most six on $\SL_2/2I$, in every degree and over
all commutative $\C$-algebras. Any remaining failure needs at least five
or seven weights, respectively, after every allowed automorphism and
choice of maximal left torus. Its all-support capacity inequality forces
one charged side to have small capacity relative to support size.
The active Mathieu thread is now exploring that inequality under
automorphisms; the intrinsic extension remains the full closure target.

The Wright notes PF126 and PF127 constrain actual peripheral permutations
and their cyclic orbits. They do not assert that the whole monodromy group
is cyclic. The following deductions add global constraints, under explicit
hypotheses, to those local restrictions.

\section{Global meridional index and genus constraints}

Let $B\subset\A^2_\C$ be a reduced branch curve. Choose a projection such
that $B$ is finite over the projection line, and let a generic transverse
fibre meet $B$ in $w$ distinct points. Assume that its geometric meridians
generate the affine complement group. This is the usual affine
Zariski--van Kamp situation. For a connected degree-$d$ finite cover,
write their permutations as $\sigma_1,\ldots,\sigma_w$.

For $\sigma\in S_d$ let $c(\sigma)$ count all cycles, including fixed points,
and put $\ind(\sigma)=d-c(\sigma)$.
\begin{proposition}
Transitivity requires
\[
             \sum_{i=1}^w\ind(\sigma_i)\ge d-1.                 \tag{1}
\]
If component $B_j$ supplies $w_j$ meridians with index $r_j$, then
$\sum_j w_jr_j\ge d-1$.
\end{proposition}
\begin{proof}
A cycle of length $a$ is a product of $a-1$ transpositions. Decompose every
$\sigma_i$ in this way. Make a graph on the $d$ sheets with these
transpositions as edges. Each original generator preserves each graph
component. A transitive group therefore forces a connected graph and at
least $d-1$ edges.
\end{proof}

The genus version retains the product. Set
$\sigma_\infty=(\sigma_1\cdots\sigma_w)^{-1}$ and
$c_\infty=c(\sigma_\infty)$. If the tuple is transitive, the associated
connected branched curve cover of $\mathbb P^1$ has genus
\[
 2g=\sum_i\ind(\sigma_i)-d-c_\infty+2.                         \tag{2}
\]
Consequently the right side is a nonnegative even integer, and
\[
       \sum_i\ind(\sigma_i)\ge d+c_\infty-2.                 \tag{3}
\]
This is classical Riemann--Hurwitz/transitive-factorization theory
\cite{gouldenjackson}, not a new general theorem. Without transitivity the
corresponding disconnected Euler-characteristic identity has a different
component term.

For transpositions, (1) gives $w\ge d-1$; for $3$-cycles it gives
$2w\ge d-1$. For example, three $3$-cycles on six sheets, if transitive,
must have $c_\infty=2$ and $g=0$ by (2). This is a necessary condition,
not exclusion of the $A_6$ case.

The public, unreviewed \href{https://github.com/alok/jacobian-two/blob/main/docs/s6-trace-curve-topology.md}{Keller trace-curve notes} give a concrete
motivation: their explicit six-sheet $S_6$ hostile trace curve has a
four-strand projection and transposition inertia, so (1) excludes its
connected cover before trace or multiplication data are considered.
A separate width-five example is reported to have cyclic complement.
Those reports eliminate their specified curves, not all degree-six
Keller candidates. The independent general argument above does not rely
on replaying their Sage computations.

\section{Affine directions give a stronger dimension bound}

Let $V=\F_2^k$ and suppose all chosen meridian images have the form
\[
 T_i(x)=x+(\ell_i(x)+a_i)u_i,\quad
 0\ne u_i\in V,\quad 0\ne\ell_i\in V^\vee,\quad\ell_i(u_i)=0.
                                                               \tag{4}
\]
These affine transvections have $d/2$ fixed points and $d/4$ two-cycles,
where $d=2^k$, so their permutation index is $d/4$.

\begin{proposition}
If the group generated by (4) is transitive on $V$, then
\[
              \langle u_1,\ldots,u_w\rangle=V,
              \qquad w\ge k.                                  \tag{5}
\]
\end{proposition}
\begin{proof}
Every generator changes a vector by an element of
$U=\langle u_1,\ldots,u_w\rangle$. Each coset of $U$ is invariant under
all words in the generators. Transitivity forces $U=V$.
\end{proof}

Combining (1) and (5) gives
\[
         w\ge\max\{k,\lceil4-4/2^k\rceil\}.                     \tag{6}
\]
This uses a single set of generators of the \emph{whole} complement group,
not just the peripheral subgroup of a knot. It cannot be applied from
local longitude data alone.

\section{A common symplectic form forces an extra meridian}

Assume in addition that a fixed nondegenerate alternating form $B$ is
preserved by all linear parts of (4). Then $k$ is even and, over $\F_2$,
\[
              \ell_i(x)=B(u_i,x).
\]
This is the same common-form hypothesis used in PF126, now coupled to
global generation.

\begin{proposition}
Under these hypotheses, transitivity forces
\[
                         w\ge k+1.                            \tag{7}
\]
More generally, if there is a quadratic refinement $q$ of $B$ with
\[
                         q(u_i)=1+a_i\quad(1\le i\le w),        \tag{8}
\]
the meridian group is not transitive.
\end{proposition}
\begin{proof}
A quadratic refinement has $q(0)=0$ and
$q(x+y)+q(x)+q(y)=B(x,y)$. Write $b=B(u_i,x)$ and
$t=b+a_i\in\F_2$. Then
\[
 q(T_i x)+q(x)
   =t\bigl(q(u_i)+b\bigr)
   =(b+a_i)(1+a_i+b)=0
\]
whenever (8) holds. Thus every $T_i$ preserves $q$. Nondegeneracy of $B$
implies that $q$ is nonconstant, so its two level sets preclude transitivity.

If $w<k$, (5) applies. If $w=k$ and the directions span, they are a basis.
Quadratic refinements form an affine space over $V^\vee$, so their values
can be prescribed arbitrarily on that basis. In particular (8) has a
solution. This rules out equality in (5).
\end{proof}

There is an inexpensive exact test for arbitrary $w$. Fix any refinement
$q_0$ and seek $L\in V^\vee$ satisfying
\[
           L(u_i)=1+a_i+q_0(u_i).                              \tag{9}
\]
A solution produces the invariant $q=q_0+L$ and excludes transitivity.
Therefore actual transitivity requires an inconsistent system (9).
Equivalently there must be a dependency $\sum_i e_i u_i=0$ with
\[
 \sum_i e_i(1+a_i)+
 \sum_{i<j}e_i e_j B(u_i,u_j)=1.                               \tag{10}
\]
This follows by applying $q_0$ to the dependency. It is independent of the
chosen affine origin: shifting the origin changes $a_i$ by $B(u_i,z)$,
whose sum in a dependency is zero.

For degree $16$, the index bound alone allows width four, whereas (7)
requires at least five in the common-symplectic branch. For degree $64$
it requires at least seven. An upper geometric width bound, or a
geometrically forced compatible refinement in (9), would turn these
necessary conditions into exclusions. Neither input has been proved here.

The potential connection is to the \emph{whole marked source} and its
trace or framing data: do they force compatible quadratic values for all
global meridians? PF126's local Seifert-framing cancellation does not by
itself provide that global refinement. Even a positive answer would close
this branch, not every remaining Keller monodromy family.

\section{Why tame coordinate reduction is not yet a Mathieu route}

Bisi--Furter--Lamy prove elementary degree reduction for nonlinear
\emph{tame automorphisms} of $\SL_2$ \cite{bfl}. Their square-complex
vertices of component type are functions completable to such an
automorphism. Neither statement covers every regular function.

There is a stronger scope obstruction in the present problem.
\begin{proposition}
Let $1\ne\Gamma\subset\SL_2(\C)$ be finite, acting by right multiplication.
No $\Gamma$-invariant regular function on $\SL_2$ can equal
$x_{11}\circ\phi$ for any algebraic automorphism $\phi$ of $\SL_2$.
\end{proposition}
\begin{proof}
For $a\ne0$, the fibre $x_{11}=a$ is $\A^2$: the remaining entries
$x_{12},x_{21}$ are arbitrary and
$x_{22}=(1+x_{12}x_{21})/a$. Thus the fibre of $x_{11}\circ\phi$ over $a$
has compactly supported Euler characteristic one. If the function were
$\Gamma$-invariant, the free right action of $\Gamma$ on $\SL_2$ would
preserve that fibre. Its quotient is a finite etale cover of degree
$|\Gamma|$, so multiplicativity gives
$1=|\Gamma|\chi_c(\text{quotient})$, impossible. Multiplicativity can
equally be proved from a finite triangulation of the covering
\cite{euler}.
\end{proof}

Thus our original exceptional phases lie outside that coordinate-component
input class, even allowing wild automorphisms. The existing strict-leading
condition after \emph{every} Haar-preserving automorphism is still strong,
but a new theorem for arbitrary invariant functions would be required to
turn it into a terminating reduction. Its proof cannot be replaced by the
published coordinate theorem. This does not disprove the possible
function-reduction strategy.

For every invariant phase and every fibre of its pullback,
$\chi_c(\text{lifted fibre})$ is divisible by $|\Gamma|$. This observation
can reject unsuitable coordinate models, but is not a Mathieu purity
criterion.

\section{The Gaussian family lead and the resonance trap}

Walther's algorithm stratifies the base of a holonomic algebraic
$D$-module so that the cohomology of each derived restriction is a
connection \cite{walther}. Interpreting these ranks as the holomorphic
solution complex uses regular holonomicity. Oaku--Takayama supply
computable restriction and local-cohomology operations
\cite{oaku}. These are relevant primary sources for the current Gaussian
family task; neither abstract identifies a specialized Gaussian polynomial
Hom space with those ranks.

For $P(c,x)=\sum_\alpha c_\alpha x^\alpha$, the joint exponential source
with exponent $F=-\sum x_i^2/2+tP$ has relations
\[
 \partial_{x_i}+x_i-tP_{x_i},\qquad
 \partial_t-P,\qquad
 \partial_{c_\alpha}-t x^\alpha.                              \tag{11}
\]
The last relations are essential. Formally,
\[
 \partial_{c_\alpha}\mathbb E[Qe^{tP}]
       =t\,\mathbb E[x^\alpha Qe^{tP}].
\]
They retain all amplitudes without justifying division by $t$ at the
origin.

\begin{lemma}
The central-parameter family $t\partial_t-c$ cannot be extended to a
rank-one flat algebraic connection on a dense open subset of
$\mathbb G_m\times\A^1_c$ while keeping its indicated $t$-connection.
\end{lemma}
\begin{proof}
Such a connection would have $\partial_t y=(c/t)y$ and
$\partial_c y=A(t,c)y$ with $A\in\C(t,c)$. Flatness implies
$\partial_t A=1/t$. A derivative of a rational function has zero residue
at every finite pole, whereas $1/t$ has residue one at zero. Contradiction.
The same trace argument works for a scalar rank-$r$ connection.
\end{proof}

Accordingly, the infinitely many integer resonances in that central
parameter example do not automatically refute a strategy using the
\emph{joint} holonomic source (11). Conversely, holonomicity does not by
itself establish all the required base-change identifications.

A complete family producer still needs specialization-valid original
amplitude coverage, the correct ordinary or derived restriction, and
finite actual-jet selection and normalization $1$ on every stratum.
Restriction to a nonlinear stratum must be proved to compute the intended
family before connection ranks or local cohomology are used. This is
aligned with the current Gaussian thread, rather than a replacement goal.

\section{Other sources checked: useful boundaries rather than closure}

Franks--Walter's positive-definite matrix-coefficient result
\cite{fw} exploits
$\int_K\langle v,\rho(k)v\rangle^m=\|\Pi_m v^{\otimes m}\|^2$.
It does not cover arbitrary complex mixed-isotypic phases, where the
positivity identity is absent. The already closed single-isotypic sector
should not be re-proposed as the remaining full problem.

Gavrilov--Pakovich \cite{gp} reduce moments on curves to inverse-branch
cycles and use monodromy. Their methods were already partially adopted
in the seed/period audits. A new application must identify the actual
boundary curve map, cycle and marker, not just a matching permutation
group.

Payne's boundary-complex description \cite{payne} controls weight-graded
cohomology. It does not determine the actual mixed extension
$\epsilon_f$. A dual-complex calculation alone cannot supply the current
Mathieu lifting step.

\href{https://github.com/long-mathematics/mathieu-second-conjecture-spherical-pairs}{Long's spherical-pair classification} was rechecked from the public repository. $\SL_2/\Gamma$ for finite $\Gamma$ has dimension three, whereas
a Borel has dimension two, so these exceptional quotients are not
spherical. His root-doublet construction remains useful within its
hypotheses; the quoted classification does not close these quotients.

\section{Priority after this search}

The most concrete broad closure target remains the Gaussian
specialization-valid family producer, now with a directly relevant
published stratification algorithm. The Mathieu thread is exploring the all-support capacity inequality under
automorphisms, following its four/six-weight closure. This is compatible
with the intrinsic extension target and does not repeat model changes.
For the Wright branch, the new coupling is (9)--(10): test whether genuine
global source geometry forces a common invariant quadratic refinement.
A purely local framing calculation or another hostile quadratic trace
matrix is insufficient.

These are ranked leads and proved necessary conditions, not solved
conjectures. The scripts accompanying this note check small permutation
and polynomial identities only; the all-degree claims rest on the
arguments above.

\begin{thebibliography}{10}
\bibitem{bfl} C. Bisi, J.-P. Furter, S. Lamy,
The tame automorphism group of an affine quadric threefold acting on a
square complex, J. Ecole polytechnique Math. 1 (2014), 161--223.
\url{https://numdam.org/articles/10.5802/jep.8/}.
\bibitem{euler} S. Cappell, L. Maxim, J. Shaneson,
Euler characteristics of algebraic varieties.
\url{https://arxiv.org/abs/math/0606654}.
\bibitem{gouldenjackson} I. P. Goulden, D. M. Jackson,
Transitive factorizations in the symmetric group, and combinatorial
aspects of singularity theory, European J. Combin. 21 (2000), 1001--1016.
\url{https://uwaterloo.ca/math/sites/default/files/uploads/documents/gjeurjc2000.pdf}.
\bibitem{walther} U. Walther,
Algorithmic stratification of RHom for regular algebraic D-modules.
\url{https://www.math.purdue.edu/~walther/research/reg.pdf}.
\bibitem{oaku} T. Oaku, N. Takayama,
Algorithms for D-modules---restriction, tensor product, localization,
and local cohomology groups.
\url{https://doi.org/10.1016/S0022-4049(00)00004-9}.
\bibitem{fw} C. Franks, M. Walter,
Minimal length in an orbit closure as a semiclassical limit.
\url{https://arxiv.org/abs/2004.14872}.
\bibitem{gp} L. Gavrilov, F. Pakovich,
Moments on Riemann surfaces and hyperelliptic Abelian integrals.
\url{https://arxiv.org/abs/1107.3029}.
\bibitem{payne} S. Payne, Boundary complexes and weight filtrations.
\url{https://arxiv.org/abs/1109.4286}.
\end{thebibliography}
\end{document}
